\documentclass[journal]{IEEEtran}

\usepackage{amsmath,amssymb}
\usepackage{booktabs}
\usepackage{graphicx}
\usepackage{cite}
\usepackage{url}
\usepackage{siunitx}
\usepackage{multirow}

\title{Physics-Gated Reliability-Constrained Adaptive Phase-Coded FMCW ISAC for High-Mobility Vehicular Links}

\author{Anonymous Author(s)%
\thanks{Manuscript prepared from the frozen publication-v2.1 simulation study and reviewer-grade supplemental evidence. Author names and affiliations are intentionally left for submission-time insertion.}}

\begin{document}
\maketitle

\begin{abstract}
This work studies adaptive phase-coded frequency-modulated continuous-wave (PC-FMCW) integrated sensing and communication (ISAC) for high-mobility vehicular links under Doppler, residual synchronization error, interference, phase noise, and imperfect PHY-state knowledge. Unlike fixed-configuration PC-FMCW feasibility studies, the proposed controller first applies a deterministic physical feasibility gate derived from waveform sampling, range, and unambiguous-velocity limits, and then performs reliability-constrained joint PHY adaptation over profile, phase-code budget, transmit-power backoff, and repetition resources. In the frozen publication-v2.1 paired benchmark, 1000 seeds and 12 scenarios per seed yield 12,000 state/seed units per policy and 72,000 receiver-level evaluations overall. The robust B4 policy selects 12.23\% of cases and satisfies joint QoS in all 1,468 selected cases under the frozen simulation protocol, corresponding to a one-sided 95\% Wilson lower bound of 99.816\%. This selected-case reliability evidence is obtained at the cost of 87.77\% abstention: B3 selects 19.07\%, and the paired B4-minus-B3 difference in unconditional joint-QoS probability is -0.03942 with a 95\% bootstrap interval [-0.04292,-0.03600]. Thus, the evidence does not support unconditional B4 superiority. Supplemental finite-draw calibration shows that a 0.95 reliability target is structurally unattainable at 32 draws; at the frozen 512-draw setting, feasibility disagreement relative to a nested 4096-draw reference is zero while selected-action disagreement remains 0.833\%. The unoptimized Python/Linux host implementation has median B4 decision latency of approximately 251.04~ms at 512 draws, with a 95th percentile of approximately 504.11~ms, so no embedded real-time claim is made. Across uncertainty scaling, ablations, distribution shift, and model mismatch, the results characterize B4 as a conservative confidence-qualified operating-region selector with an explicit reliability--availability--computation trade-off rather than as a universally superior policy.
\end{abstract}


\begin{IEEEkeywords}
Integrated sensing and communication, phase-coded FMCW, automotive radar, robust PHY adaptation, chance constraints, high mobility, physical feasibility, vehicular communications.
\end{IEEEkeywords}

\section{Introduction}
\label{sec:introduction}
Integrated sensing and communication (ISAC) seeks to share waveform, spectrum, and hardware resources between sensing and communication while exposing explicit performance trade-offs between the two functions \cite{liu2022isac,wang2022isac}. Phase-coded frequency-modulated continuous-wave (PC-FMCW) signaling is attractive in this setting because a common FMCW carrier can retain range--Doppler sensing structure while phase coding carries communication information. Prior PC-FMCW studies have demonstrated the sensing--communication principle experimentally, analyzed receiver-dependent sensing--communication trade-offs, and investigated interference resilience \cite{kumbul2021experimental,kumbul2023performance,kumbul2024interference}. The waveform family has also been developed through dedicated receiver-structure studies, smoothed phase-coded transceiver architectures, coherent MIMO operation, and code-dependent sensing analysis \cite{kumbul2022receiver,kumbul2022smoothed,kumbul2023mimo,kumbul2022codes}. More recently, phase-coded FMCW has been used in an optical laser-headlamp architecture that combines sensing, communication, and illumination \cite{liu2026laserheadlamp}. These results establish waveform and functional feasibility, but they do not resolve a different problem that appears in high-mobility vehicular operation: a configuration that is nominally attractive can be physically incapable of representing the requested range or radial velocity, or can violate communication and sensing QoS once state uncertainty, synchronization residuals, and interference are considered.

High mobility also makes reliability under uncertain link state a first-order concern. Robust V2X resource-allocation work has shown that imperfect channel knowledge and heterogeneous QoS requirements motivate explicit chance-constrained or distributionally robust treatment rather than deterministic optimization alone \cite{wu2021robustv2x}. Recent robust ISAC work similarly treats imperfect communication CSI and uncertain target location as coupled design uncertainties \cite{lyu2024dualrobust}. A 2026 survey of vehicular ISAC identifies robustness to imperfect modeling and hardware non-idealities, scalable real-time control, and prototype/field validation among the key obstacles to practical deployment \cite{ma2026vehicularisac}. The present problem differs in control variable and evidence target: instead of optimizing a continuous beamformer or allocating network resources across users, we adapt a finite PC-FMCW waveform/PHY action while simultaneously enforcing deterministic waveform feasibility and joint sensing--communication reliability.

This paper therefore moves the research question from fixed-configuration feasibility to reliable adaptive operation. We study an RF/mmWave \SI{77}{GHz} vehicular PC-FMCW physical layer and ask which configuration should be selected, given imperfect knowledge of the link and sensing state, so that communication and sensing QoS are jointly supported with a declared reliability while physical resources remain controlled. The design uses two ordered stages. A PC-FMCW-specific physical feasibility gate first removes configurations that violate sampling, range, or unambiguous-velocity limits. A robust reliability selector then chooses the minimum-cost admissible action satisfying joint communication and sensing constraints under an uncertainty distribution.

The main contributions are as follows:
\begin{itemize}
    \item A reproducible PC-FMCW system model that explicitly separates the one-way remote communication path from the two-way monostatic sensing echo and models radar sampling at dechirped IF rather than at the \SI{77}{GHz} carrier.
    \item A physics-gated finite-action adaptation framework based on waveform bandwidth, chirp slope, IF sampling, maximum range, chirp repetition interval, and unambiguous radial velocity.
    \item A reliability-constrained joint PHY selector over waveform profile, phase-code/chip budget, transmit-power backoff, and repetition resources under uncertain SNR, residual CFO/synchronization error, interference, and phase-noise state.
    \item A finite-sample confidence rule in which an action is accepted only when a one-sided Wilson lower confidence bound reaches the declared joint-reliability target, together with calibration of the statistical and computational consequences of the draw budget.
    \item A feasible-operating-region evaluation that distinguishes hard physical infeasibility, policy abstention, and selected-action failure, rather than collapsing these outcomes into a single success rate.
    \item A frozen paired Monte-Carlo comparison and reviewer-grade uncertainty, ablation, mismatch, Pareto, distribution-shift, and runtime analysis demonstrating a reliability--availability--complexity trade-off instead of unconditional superiority.
\end{itemize}

The scope is deliberately PHY-level. The work does not perform trajectory forecasting, ego-motion planning, packet/user scheduling, beam management, adaptive driving-beam illumination, or optical headlamp control.

\section{Related Work and Positioning}
\label{sec:related}
Broad ISAC surveys identify waveform design and the sensing--communication performance trade-off as central physical-layer design questions, with vehicular/V2X operation among the relevant application domains \cite{liu2022isac,wang2022isac}. Within the PC-FMCW family, experimental work established that phase coding can combine communication capability with FMCW-style sensing \cite{kumbul2021experimental}. Receiver-focused studies then analyzed practical low-sampling receiver structures and sensing/communication behavior, including the dependence of communication BER and sensing performance on the phase-modulation processing choice \cite{kumbul2022receiver,kumbul2023performance}. Smoothed PC-FMCW work further developed the waveform/transceiver architecture and evaluated it through simulation and experiments \cite{kumbul2022smoothed}. Related PC-FMCW research addresses coherent MIMO operation and code-dependent sensing behavior \cite{kumbul2023mimo,kumbul2022codes}, while separate work studies mitigation of mutual automotive-radar interference \cite{kumbul2024interference}. These contributions motivate the waveform family, sensing model, and impairment axes used here, but the present paper does not claim novelty for phase coding itself, DPSK embedding, range--Doppler processing, MIMO coding, receiver filtering, or generic joint sensing and communication.

Robust optimization under uncertain state is also established. Distributionally robust V2X resource allocation has been formulated around chance constraints to protect communication QoS under imperfect CSI \cite{wu2021robustv2x}; robust ISAC beamforming has considered CSI imperfection jointly with target-location uncertainty \cite{lyu2024dualrobust}. More generally, the scenario approach connects sampled constraints with chance-constrained feasibility \cite{campi2011sampling,care2014fast}. Our use of random draws and confidence qualification is not presented as a new theory of chance constraints. The distinction is the PC-FMCW-specific composition of (i) a hard physical capability gate, (ii) uncertain vehicular PHY state including synchronization and interference impairments, (iii) joint communication and sensing QoS, (iv) a finite discrete action set, and (v) explicit abstention when the available evidence is insufficient.

The abstention mechanism also has a useful conceptual analogue in selective prediction: rejecting uncertain cases can reduce risk on the covered domain at the cost of lower coverage \cite{geifman2019selectivenet}. Here, selection rate plays the role of operational coverage. Consequently, the manuscript always reports conditional reliability together with unconditional joint-QoS probability and abstention; a high selected-case success rate alone is not interpreted as universal service reliability.

Vehicular ISAC is moving toward standardization and deployment-oriented study. 3GPP TR 22.837 is a Release-19 study on integrated sensing and communication \cite{3gpp22837}, while recent vehicular-ISAC literature emphasizes high mobility, robustness, scalability, and hardware/model non-idealities \cite{ma2026vehicularisac}. These developments motivate the deployment context but do not convert the present simulation study into a standards-compliant implementation or hardware validation.

The closest conceptual predecessor to the broader vehicular integration theme is the phase-coded FMCW laser-headlamp ISCAI architecture \cite{liu2026laserheadlamp}. That system addresses optical sensing, communication, and adaptive illumination. The present work is independent in both medium and optimization objective: it studies a \SI{77}{GHz} RF/mmWave PC-FMCW PHY and adapts the physical-layer configuration under high-mobility uncertainty. The unit of control is neither the vehicle trajectory nor a packet queue, beam, or illumination pattern.

The distinguishing question is therefore not whether PC-FMCW can support sensing and communication, nor whether robust optimization can be used in isolation, but where a physically admissible PC-FMCW configuration can jointly support both functions at a declared reliability under mobility-related uncertainty. This motivates a two-layer formulation in which deterministic waveform physics is enforced before stochastic reliability qualification.

\section{System Model}
\label{sec:system}
\subsection{PC-FMCW Transmit Signal}
For chirp index $m$ and fast time $t$, the idealized transmitted waveform is
\begin{equation}
 s_m(t)=\sqrt{P_m}\exp\left\{j\left[\pi\mu t^2+\phi_m(t)\right]\right\},
\end{equation}
where $\mu=B/T_c$ is the chirp slope, $B$ is sweep bandwidth, $T_c$ is active chirp duration, $P_m$ is transmit power, and $\phi_m(t)$ is the phase-coded communication sequence. Multiple phase-code chips may be carried in one chirp.

\subsection{Monostatic Sensing Path}
For target range $R$ and radial velocity $v$, the sensing echo uses the two-way delay and Doppler
\begin{equation}
 \tau_r=\frac{2R}{c}, \qquad f_{D,r}=\frac{2vf_c}{c}.
\end{equation}
After dechirping, the practical radar ADC samples the IF/beat signal, not the \SI{77}{GHz} carrier. Under the sign convention used in the reference simulator,
\begin{equation}
 f_b=\mu\tau_r+f_{D,r}.
\end{equation}
The implementation synthesizes a two-dimensional slow-time/fast-time complex IF matrix. Range--Doppler processing uses windowed FFTs; Doppler is estimated first and removed from the beat-frequency estimate before range recovery, reducing the standard FMCW range--Doppler coupling bias. Receiver-level sensing performance is evaluated through repeated noisy IF trials and RMSE of the resulting range and radial-velocity errors.

\subsection{One-Way Communication Path}
The remote communication receiver experiences
\begin{equation}
 \tau_c=\frac{R}{c}, \qquad f_{D,c}=\frac{vf_c}{c}.
\end{equation}
The synchronized receiver removes the known FMCW chirp and decodes the remaining phase sequence. A transparent multi-chip DBPSK modem is used as the communication reference. In AWGN its analytical noncoherent bit-error probability is
\begin{equation}
 P_b=\frac{1}{2}\exp(-E_b/N_0).
\end{equation}
For $N_{\rm chip}$ chips per chirp and chirp repetition interval $T_r$, the reference raw information rate is $N_{\rm chip}/T_r$. Residual frequency error rotates the baseband chip sequence according to the physical chip times, including the inter-chirp idle interval. The final receiver evaluation additionally includes the declared interference and phase-noise impairments and applies repetition combining according to the selected action.

\section{Physical Feasibility and Adaptive PHY Formulation}
\label{sec:method}
\subsection{PC-FMCW Capability Gate}
For each profile, the simulator derives the fundamental capability quantities
\begin{align}
 \Delta R &= \frac{c}{2B},\\
 R_{\max}^{(+\mathrm{IF})} &= \frac{cF_s}{4\mu},\\
 v_{\max} &= \frac{\lambda}{4T_r},\\
 \Delta v &= \frac{\lambda}{2N_cT_r},
\end{align}
where $F_s$ is IF sample rate and $N_c$ is the number of chirps used for slow-time processing. A candidate profile is excluded whenever the requested operating state lies outside its declared physical support. This gate is deterministic and precedes any empirical or Monte-Carlo reliability estimate.

The source-grounded short-range profile uses a \SI{77}{GHz} carrier, \SI{858}{MHz} valid sweep, $T_c=\SI{25.6}{\micro s}$, $T_r=\SI{115.8}{\micro s}$, $F_s=\SI{10}{MSPS}$, 256 samples/chirp, and 64 chirps/frame, following a Texas Instruments automated-parking reference profile \cite{tiParking}. Its range resolution is approximately \SI{0.175}{m}, positive-IF range support approximately \SI{22.36}{m}, radial-velocity resolution approximately \SI{0.263}{m/s}, and unambiguous radial-velocity support approximately \SI{8.41}{m/s}. A separate high-mobility capability profile uses a \SI{1}{GHz} sweep, \SI{20}{\micro s} active/repetition timing, \SI{37.5}{MSPS} sampling, 750 samples/chirp, and 128 chirps/frame, giving approximately \SI{0.150}{m} range resolution, \SI{56.21}{m} positive-IF range support, \SI{0.760}{m/s} velocity resolution, and \SI{48.67}{m/s} unambiguous radial-velocity support. The latter is a composite capability reference, not a claimed commercial preset.

\subsection{Finite PHY Action Space}
The frozen action set is the Cartesian product of two profiles, three chip budgets $\{16,32,64\}$, three transmit-power backoffs $\{0,3,6\}$ dB, and three repetition factors $\{1,2,4\}$, yielding
\begin{equation}
 |\mathcal{A}|=2\times3\times3\times3=54.
\end{equation}
The dimensions are intentionally interpretable rather than latent optimization variables: profile choice changes physical range/velocity and sampling burden; chips per chirp changes raw communication rate and chip duration; transmit-power backoff changes communication and sensing SNR; and repetition changes rate and error protection.

The declared normalized resource cost is a dimensionless experimental index rather than energy in joules. For action $a$ it is implemented as
\begin{equation}
 C(a)=g_p\left(0.55\,10^{-b/10}+0.25\frac{N_{\rm chip}}{16}+0.20N_{\rm rep}\right),
\end{equation}
where $b$ is power backoff in dB, $N_{\rm rep}$ is the repetition factor, and $g_p=1$ for the parking profile and $1.6$ for the high-mobility capability profile. This synthetic index is used only for consistent action ordering and must not be interpreted as physical energy consumption.

Let $S$ denote the uncertain physical/link state, including $E_b/N_0$, IF SNR, range, radial velocity, residual CFO, interference, phase-noise state, and estimation uncertainty. The physical gate defines $\mathcal{A}_{\mathrm{phys}}(\hat S)\subseteq\mathcal{A}$. Among admissible actions, the robust controller seeks the minimum-cost configuration satisfying the declared joint-QoS reliability.

\subsection{Baselines and Oracle}
B0 uses a fixed high-mobility configuration whenever that configuration remains physically admissible. B1 minimizes resource cost subject to communication-only selector constraints, B2 does the analogous sensing-only adaptation, and B3 requires communication and sensing constraints jointly while treating the estimated state as exact. B4 adds uncertainty draws and confidence qualification. The Oracle is deliberately non-deployable: it uses the true state and receiver-level evaluation to find the minimum declared resource-cost action that succeeds in hindsight. The Oracle is therefore a reference bound, not a realizable controller.

\section{Finite-Sample Reliability Qualification}
\label{sec:confidence}
For a physically admissible action $a$, B4 draws uncertain states around the estimated state and records the binary joint-QoS outcome. If $k$ of $N$ draws succeed, $\hat p=k/N$ is an empirical fraction, not a reliability certificate. The implementation therefore uses a one-sided Wilson lower confidence bound
\begin{equation}
 L_W=\frac{\hat p+\frac{z^2}{2N}-z\sqrt{\frac{\hat p(1-\hat p)}{N}+\frac{z^2}{4N^2}}}{1+\frac{z^2}{N}},
\end{equation}
where $z=\Phi^{-1}(0.95)\approx1.64485$ for the frozen one-sided 95\% confidence rule. Wilson-type intervals avoid the well-known poor small-sample/boundary behavior of the elementary Wald interval \cite{brown2001binomial}. B4 admits an action only if
\begin{equation}
 L_W\ge\rho,\qquad \rho=0.95,
\end{equation}
and then selects the cheapest admitted action. If no action qualifies, the controller abstains.

The finite-draw budget has a structural effect. With all $N$ trials successful, the Wilson lower bound simplifies to $N/(N+z^2)$. Thus, with $N=32$, even 32/32 successes give $L_W\approx0.9220$ and the frozen target is unattainable. At 64 draws the maximum lower bound is approximately 0.9594, so the target becomes attainable, but only with 64/64 successes. The supplemental calibration later evaluates 32, 64, 128, 256, and 512 draws against a nested 4096-draw reference. That high-draw reference reduces Monte-Carlo decision noise; it is not physical ground truth.

B4 uses an action-local deterministic uncertainty stream derived from the state seed and action identity. This prevents the robust-selector sampling stream from being conflated with the held-out final receiver evaluation. The frozen protocol uses 512 internal draws and does not tune the reliability threshold after observing final results.

\section{Evaluation Protocol}
\label{sec:protocol}
The frozen publication-v2.1 protocol declares communication BER $\le10^{-3}$, effective rate $\ge\SI{100}{kbps}$, range RMSE $\le\SI{1}{m}$, velocity RMSE $\le\SI{1}{m/s}$, and a target joint reliability of 0.95. The v2.1 change corrects a selector-side sensing-SNR scaling defect identified by the mandatory smoke gate; the QoS thresholds, receiver-level evaluation, action set, seed bank, and statistical rules were left unchanged.

Six policies are evaluated on common support: B0, B1, B2, B3, B4, and Oracle. The frozen primary benchmark uses 12 predeclared scenarios for each of 1000 seeds beginning at 10000, yielding 12000 state/seed units per policy and 72000 receiver-level evaluations overall. Each final receiver evaluation uses 20000 communication bits and three sensing trials. Policy comparisons are paired on the common state/seed bank. Paired uncertainty intervals use 10000 bootstrap resamples at 95\% confidence. This paired design removes between-bank Monte-Carlo variation from policy differences and is therefore more informative than comparing independent random test sets.

Primary results are frozen separately from reviewer-grade supplemental evidence. Smoke runs are structural diagnostics and are not promoted to publication results. Supplemental studies probe uncertainty scale, reliability targets, finite-draw calibration, physics and uncertainty-source ablations, distribution shift, model mismatch, action-space restrictions, confidence maps, empirical Pareto structure, and host runtime.

All numerical outputs in this paper are controlled simulation, analytical, statistical, or host-runtime evidence. Literature-grounded hardware constants and externally reported measurements are retained as provenance and are not represented as measurements produced by this work.

% Auto-generated from frozen v2.1 summary JSON. Do not edit numerical values manually.
\begin{table*}[t]
\centering
\caption{Frozen v2.1 policy-level reliability and availability results.}
\label{tab:v21_policy}
\begin{tabular}{lrrrr}
\hline
Policy & Selection (\%) & Joint QoS (\%) & Conditional QoS (\%) & 95\% Wilson LCB (cond.) \\
\hline
B0 FIXED & 92.36 & 59.04 & 63.93 & 63.17 \\
B1 COMM ONLY & 75.32 & 52.59 & 69.83 & 69.03 \\
B2 SENSING ONLY & 32.26 & 16.18 & 50.14 & 48.82 \\
B3 DETERMINISTIC JOINT & 19.07 & 16.18 & 84.83 & 83.56 \\
B4 ROBUST JOINT & 12.23 & 12.23 & 100.00 & 99.82 \\
ORACLE & 66.67 & 66.67 & 100.00 & 99.97 \\
\hline
\end{tabular}
\end{table*}

\begin{table*}[t]
\centering
\caption{Frozen v2.1 E11 ablation results.}
\label{tab:v21_e11}
\begin{tabular}{lrrr}
\hline
Variant & Selection (\%) & Joint QoS (\%) & Conditional QoS (\%) \\
\hline
FULL B4 & 12.50 & 12.50 & 100.00 \\
NO STATE UNCERTAINTY & 16.67 & 16.67 & 100.00 \\
NO CFO & 12.50 & 12.50 & 100.00 \\
NO INTERFERENCE & 12.50 & 12.50 & 100.00 \\
\hline
\end{tabular}
\end{table*}

% Frozen paired-bootstrap result for manuscript text:
% B4-B3 mean difference = -3.942 percentage points; 95% CI [-4.292, -3.600] pp.
% B4-B0 mean difference = -46.808 percentage points; 95% CI [-47.700, -45.917] pp.
% Claim boundary: these results do NOT support unconditional superiority of B4.
% Hindsight ORACLE is a non-deployable reference and must be labeled as such.

\begin{table*}[t]
\centering
\caption{Supplemental uncertainty sweep: B3 versus B4. Values are selection rate and conditional joint-QoS probability under the evaluated uncertainty multiplier.}
\label{tab:supp-uncertainty}
\begin{tabular}{c|cc|cc}
\hline
Uncertainty scale & B3 selection & B3 conditional QoS & B4 selection & B4 conditional QoS\\
\hline
0.0 & 16.67\% & 100.00\% & 16.67\% & 100.00\%\\
0.5 & 16.67\% & 100.00\% & 16.17\% & 100.00\%\\
1.0 & 18.50\% & 87.39\% & 12.50\% & 100.00\%\\
1.5 & 19.75\% & 74.68\% & 11.17\% & 100.00\%\\
2.0 & 20.33\% & 68.03\% & 9.75\% & 100.00\%\\
3.0 & 22.83\% & 62.04\% & 5.92\% & 98.59\%\\
\hline
\end{tabular}
\end{table*}

\begin{table}[t]
\centering
\caption{Reviewer-grade supplemental ablation study.}
\label{tab:supp-ablation}
\begin{tabular}{lcc}
\hline
Variant & Selection & Conditional joint QoS\\
\hline
FULL B4 & 12.50\% & 100.00\%\\
No physics gate & 31.67\% & 0.00\%\\
No state uncertainty & 20.00\% & 80.83\%\\
No joint constraint & 67.25\% & 85.50\%\\
\hline
\end{tabular}
\end{table}

\begin{table*}[t]
\centering
\caption{Finite-draw confidence calibration against the nested 4096-draw reference. Zero feasibility disagreement at 512 draws does not imply identical action selection.}
\label{tab:finite-draw-calibration}
\begin{tabular}{rccccc}
\hline
Draws & Max. Wilson LCB & Min. successes & Feasibility disagreement & False feasible & Selected-action disagreement\\
\hline
32  & 0.9220 & unattainable & 12.583\% & 0.000\% & 12.583\%\\
64  & 0.9594 & 64/64  & 0.583\% & 0.083\% & 3.583\%\\
128 & 0.9793 & 126/128 & 0.333\% & -- & 2.167\%\\
256 & 0.9895 & 249/256 & 0.083\% & -- & 1.250\%\\
512 & 0.9947 & 495/512 & 0.000\% & 0.000\% & 0.833\%\\
\hline
\end{tabular}
\end{table*}

\begin{table}[t]
\centering
\caption{Decision-time scaling of the unoptimized Python/Linux host implementation in the final reviewer-grade runtime evidence. These values are not embedded-ECU timings.}
\label{tab:runtime-scaling}
\begin{tabular}{rccc}
\hline
Robust draws & B3 median [ms] & B4 median [ms] & B4 p95 [ms]\\
\hline
64  & $\approx0.50$ & 32.17 & 64.37\\
128 & $\approx0.50$ & 63.81 & 127.98\\
256 & $\approx0.50$ & 127.19 & 255.28\\
512 & $\approx0.50$ & 251.04 & 504.11\\
\hline
\end{tabular}
\end{table}

% Final Results and Discussion narrative for publication v2.1.
% Evidence basis:
% - frozen main benchmark: GitHub Actions run 33967030983
% - supplemental reviewer evidence: GitHub Actions run 33975266575
% All numerical values are simulation/analytical outputs, not hardware measurements.

\section{Results and Discussion}
\label{sec:results-discussion}

\subsection{Frozen paired benchmark}

The frozen v2.1 benchmark evaluates six policies over the same 12,000 operating-state/seed pairs per policy, yielding 72,000 receiver-level evaluations in total. The fixed policy B0 selects in 92.36\% of cases and attains 59.04\% unconditional joint QoS, whereas the deterministic joint controller B3 selects in 19.07\% of cases and attains 16.18\% unconditional joint QoS. The proposed robust controller B4 is substantially more selective: it acts in only 12.23\% of evaluated cases and abstains in 87.77\%. Its unconditional joint-QoS probability is therefore 12.23\%, lower than both B0 and B3.

This reduction in unconditional availability is not a statistical ambiguity. The paired B4--B3 difference in unconditional joint QoS is $-0.03942$, with a 95\% paired-bootstrap interval of $[-0.04292,-0.03600]$ over 12,000 paired observations and 10,000 bootstrap resamples. Likewise, the B4--B0 difference is $-0.46808$ with a 95\% interval of $[-0.47700,-0.45917]$. Consequently, the frozen evidence does not support a claim that robust adaptation maximizes unconditional joint QoS.

Conditioned on selection, however, the behavior changes qualitatively. B3 achieves 84.83\% conditional joint QoS, while B4 succeeds in all 1,468 selected cases. The corresponding one-sided 95\% Wilson lower bound for B4 is 99.816\%, well above the declared 95\% reliability target. The central empirical result is therefore a reliability--availability trade-off: B4 does not enlarge the operating region; it contracts the operating region until the retained states are highly reliable under the declared uncertainty model.

The hindsight Oracle selects 66.67\% of cases and has conditional joint QoS equal to one by construction because it uses true receiver-level state information unavailable to deployable policies. It is therefore reported only as a reference, not as an implementable upper bound on computationally causal control.

\subsection{Effect of state uncertainty}

The supplemental uncertainty sweep isolates the distinction between deterministic and robust joint adaptation. At uncertainty scale zero, B3 and B4 coincide: both select 16.67\% of cases and achieve 100\% conditional joint QoS. As uncertainty increases, B3 continues to admit a comparatively broad set of states, but its conditional reliability degrades. At uncertainty scale one, B3 selects 18.50\% with 87.39\% conditional joint QoS, while B4 selects 12.50\% and preserves 100\%. At scale two, B3 selects 20.33\% with 68.03\% conditional reliability, whereas B4 contracts to 9.75\% selection and again retains 100\% observed conditional joint QoS.

At the most severe tested scale of three, B3 conditional joint QoS falls to 62.04\%. B4 selects only 5.92\% of cases and achieves 98.59\% observed conditional joint QoS. This point must be interpreted cautiously: its one-sided 95\% Wilson lower bound is 93.93\%, below the declared 95\% target. Thus, the strongest uncertainty setting identifies a confidence-qualified failure boundary even though the raw point estimate remains high. This is consistent with the intended role of B4: robust uncertainty treatment preserves reliability over a smaller region, but it does not provide unlimited protection as model uncertainty grows.

\subsection{Why the physics gate and joint constraint matter}

The extended ablation study separates the effects of physical feasibility, uncertainty modeling, and the joint sensing--communication constraint. FULL\_B4 selects 12.50\% of the ablation bank, achieves 100\% observed conditional joint QoS, and has a 95\% Wilson lower bound of 98.23\%. Removing the state-uncertainty model expands selection to 20.00\% but reduces conditional joint QoS to 80.83\%. Removing the joint constraint expands selection further to 67.25\%, while conditional joint QoS decreases to 85.50\%.

The most extreme ablation is the removal of the physical feasibility gate. In this frozen supplemental design, the ungated variant selects 31.67\% of states but achieves 0\% joint QoS. This does not imply that every conceivable ungated controller must fail. Rather, it demonstrates that optimization over candidate PHY profiles is scientifically ill-posed if actions that violate hard range or unambiguous-velocity limits are allowed to compete with physically admissible actions.

The source-grounded profile limits make this distinction explicit. The short-range parking profile supports approximately 22.36~m positive-IF range and 8.41~m/s unambiguous radial velocity, while the high-mobility capability profile supports approximately 56.21~m and 48.67~m/s. The resulting range--velocity map therefore contains three physically distinct regions: no feasible profile, high-mobility-only feasibility, and overlap where both profiles are admissible. Policy abstention should not be conflated with this hard physical infeasibility.

\subsection{Model mismatch and impairment robustness}

The supplemental mismatch experiments show that robust adaptation can provide meaningful resilience to moderate under-modeling, while also exposing clear failure boundaries. When the controller assumes 500~Hz residual CFO but the actual residual is between 1 and 5~kHz, B3 attains approximately 87--88\% joint QoS in the evaluated 100-seed banks, whereas B4 attains 100\%. When a 20~m/s Doppler condition is assumed and the actual radial speed is 25--40~m/s, B3 attains 86\% and B4 again attains 100\% in the tested banks. For interference under-modeling from assumed INR $-10$~dB to actual INR 0~dB, B3 falls to 18\%, while B4 reaches 86\%.

These gains are not universal. At actual INR values of 10 or 20~dB, both B3 and B4 fail in the tested mismatch bank. Likewise, when the controller assumes 12~dB SNR but the actual SNR is only 8~dB, B4 joint QoS falls to 17\%, and at 6~dB both policies fail. These negative results are important because they delimit the uncertainty set within which robust adaptation remains useful. They also prevent interpreting the proposed controller as a substitute for accurate synchronization, interference suppression, or link-quality estimation.

\subsection{Communication and sensing impairment evidence}

The waveform-level validation provides an independent physical explanation for the policy-level behavior. For differential communication at 8~dB, the residual-frequency stress test shows BER increasing from approximately $9\times10^{-4}$ near zero residual offset to approximately $3.2\times10^{-2}$ around 5~kHz. This establishes that residual Doppler/synchronization error at vehicular scales can materially affect communication reliability and therefore must remain an explicit state variable.

The sensing validation similarly shows a sharp SNR-dependent transition rather than a uniformly benign estimator. For the short-range profile at 10~m and 3~m/s, range RMSE is several metres at very low IF-SNR but falls to centimetre-scale once the estimator is within its usable regime. For the high-mobility profile at approximately 20.28~m and 19.86~m/s, the target lies outside the parking profile's unambiguous-velocity support but within the high-mobility profile's physical support. This is precisely the type of case for which physics-gated profile selection is required before statistical reliability optimization.

\subsection{Resource trade-offs}

The reported resource analysis separates physical decision variables from the experiment-defined normalized cost. The supplemental Pareto representation reports transmit-power fraction, chips per chirp, repetition factor, ADC samples per frame, effective rate, and sensing accuracy directly. In the frozen main output, one reported B4 operating point uses 16 chips/chirp, repetition factor one, transmit-power fraction 0.2512, and 96,000 ADC samples/frame, with observed conditional joint QoS equal to one over 828 observations. The normalized resource cost associated with the frozen optimization is dimensionless and must not be interpreted as energy in joules.

The broader implication is that robust reliability is not obtained for free. B4 sacrifices operating-region availability and, for some selected states, requires a more conservative resource allocation than deterministic adaptation. The correct comparison is therefore multi-objective: reliability, availability, effective rate, sensing accuracy, and physical resource usage must be considered jointly rather than collapsed into a single unconditional-success metric.

\subsection{Runtime and implementation limitation}

The current Python reference implementation is not yet a real-time vehicular controller at the frozen B4 setting. Median B3 decision latency is approximately 0.44~ms. B4 scales approximately linearly with the number of uncertainty draws: 30.71~ms at 64 draws, 61.03~ms at 128, 121.36~ms at 256, and 240.44~ms at 512. These are prototype timings on the GitHub Actions execution platform, not embedded-target measurements.

Accordingly, no real-time deployment claim is made for the 512-draw implementation. The timing result instead motivates algorithmic optimization through action pruning after the physical gate, vectorized uncertainty evaluation, parallel sampling, cached surrogate models, and reduced-draw or sequential confidence tests. The physical gate is especially useful here because it can remove impossible candidates before the expensive robust reliability evaluation.

\subsection{Overall interpretation and claim boundary}

Taken together, the frozen and supplemental evidence supports a narrower and more defensible conclusion than an unconditional-superiority narrative. Robust physics-gated PC-FMCW adaptation identifies a conservative subset of the high-mobility operating space in which selected actions achieve substantially higher conditional joint sensing--communication reliability under state uncertainty and several forms of moderate model mismatch. This gain is purchased by a pronounced reduction in availability and by higher computational cost.

The evidence does not support the claims that B4 maximizes unconditional joint QoS, remains reliable under arbitrary SNR or interference mismatch, or is already real-time at 512 uncertainty draws. All reported numerical results are simulation or analytical evidence; none are hardware measurements. The principal contribution is therefore the explicit characterization of the physically feasible and statistically reliable operating region, together with the demonstrated reliability--availability--complexity trade-off induced by robust PC-FMCW PHY adaptation.


\section{Reproducibility and Claim Boundaries}
\label{sec:reproducibility}
The public research repository freezes the publication-v2.1 protocol, fixed seed family, policy implementation, receiver models, machine-readable JSON/CSV results, table generators, supplemental evidence scripts, and CI workflows. The manuscript source is built by an IEEEtran LaTeX audit workflow that fails on unresolved citations or references. A separate CI workflow runs the unit suite and smoke validations spanning the communication/sensing validation path, the paired policy benchmark, and the supplemental research surface.

Reproducibility does not widen the claim class. The conditional reliability statements are conditional on the declared uncertainty models and finite action set. The empirical Pareto partition is restricted to realized evaluated operating points. The 4096-draw calibration reference is not physical truth. Runtime is measured for the reference Python/Linux implementation and is not an embedded-ECU timing certification. Packet error rate is not independently measured without a declared packet model. Finally, the composite high-mobility profile is a simulator capability reference rather than a statement that a particular commercial radar exposes the identical preset.

\section{Limitations}
\label{sec:limitations}
The reliability conclusions are conditional on the declared uncertainty families and finite action set. Severe SNR or interference mismatch can invalidate the robust operating region, as shown in the supplemental mismatch banks. The communication receiver is a transparent DBPSK reference implementation rather than a measured automotive communication standard, and the phase-noise specification is not converted into a detailed oscillator model without additional assumptions. The high-mobility profile is a composite capability reference. Finally, runtime is measured for an unoptimized Python implementation rather than an embedded automotive target; hardware-in-the-loop and RF experiments are required before deployment-level latency, energy, or measurement claims can be made.

\section{Conclusion}
\label{sec:conclusion}
This paper extends PC-FMCW ISAC from fixed-configuration functional feasibility toward physics-aware and reliability-constrained operation under high-mobility uncertainty. The proposed methodology separates whether a candidate waveform/profile is physically capable of supporting the requested range and radial velocity from whether the remaining admissible candidates satisfy joint communication and sensing reliability under uncertain PHY conditions. The physical gate prevents unsupported profiles from entering statistical selection, while the robust controller explicitly trades operating-region availability for stronger selected-case reliability evidence.

The frozen paired benchmark does not show unconditional superiority of the robust controller. B4 selects 12.23\% of evaluated cases and succeeds in all 1,468 selected cases, with a one-sided 95\% Wilson lower bound of 99.816\%. B3 selects 19.07\% and has higher unconditional joint-QoS probability. The frozen paired B4-minus-B3 unconditional effect is $-0.03942$, with a 95\% bootstrap interval of $[-0.04292,-0.03600]$. The independent 1,200-pair supplemental full-metric bank reproduces the same qualitative trade-off: B4 selects 12.50\% with 100\% observed conditional joint QoS and a 98.23\% Wilson lower bound, whereas B3 selects 20.00\% with 80.83\% conditional joint QoS; the paired unconditional effect is $-0.03667$ with 95\% interval $[-0.04833,-0.02583]$. Universal or unconditional B4 superiority is therefore explicitly rejected.

Ablation and model-mismatch experiments further show that physical feasibility, uncertainty modeling, and the joint constraint are consequential. Moderate evaluated CFO, Doppler, and interference mismatch can favor B4 on the tested banks, but severe SNR and interference mismatch remain failure modes. Finite-draw calibration against a nested 4096-draw reference shows that 32 draws cannot attain the declared 0.95 reliability target under the one-sided Wilson acceptance rule. At the frozen 512-draw setting, feasibility-decision disagreement is zero relative to that reference, while selected-action disagreement remains 0.833\%; the high-draw comparator is not physical ground truth.

The present implementation also exposes a computational limitation. B4 median host decision time is approximately 32.17, 63.81, 127.19, and 251.04~ms for 64, 128, 256, and 512 uncertainty draws, respectively; at 512 draws the 95th percentile is approximately 504.11~ms. These are unoptimized Python/Linux host measurements, not embedded-ECU certification. Hardware validation and algorithmic acceleration are therefore necessary before deployment-level latency, energy, or measurement claims can be made.

Overall, the central result is not that robust adaptation wins every operating point. Rather, physics-gated robust PC-FMCW adaptation makes the reliability boundary explicit: it identifies evaluated operating states in which joint sensing--communication service satisfies the adopted confidence-qualified simulation criterion, states in which the controller should abstain, and states in which the underlying waveform physics makes the requested service infeasible. The contribution is thus a reliability--availability--computation characterization of a finite PC-FMCW operating region, not a claim of global Pareto optimality or universal policy dominance.


\bibliographystyle{IEEEtran}
\bibliography{references_v2_1,references_submission}

\end{document}
